\chapter*{Danh mục thuật ngữ Anh--Việt}
\addcontentsline{toc}{chapter}{Danh mục thuật ngữ Anh--Việt}
\begingroup
\small
\setstretch{1.0}
\setlength{\tabcolsep}{5pt}
\renewcommand{\arraystretch}{1.15}
\begin{longtable}{@{}p{0.33\textwidth}p{0.63\textwidth}@{}}
\toprule
\textbf{Thuật ngữ tiếng Anh} & \textbf{Nghĩa và cách dùng trong khóa luận} \\
\midrule
\endfirsthead
\multicolumn{2}{l}{\textit{Danh mục thuật ngữ Anh--Việt (tiếp theo)}}\\
\toprule
\textbf{Thuật ngữ tiếng Anh} & \textbf{Nghĩa và cách dùng trong khóa luận} \\
\midrule
\endhead
\bottomrule
\endlastfoot
Ablation study & Nghiên cứu loại bỏ hoặc thay thế thành phần để khảo sát ảnh hưởng của lựa chọn thiết kế. \\
AUROC-OOD & Diện tích dưới đường cong ROC dùng để đánh giá khả năng xếp hạng mẫu ID và OOD. \\
AUPR-Out & Diện tích dưới đường cong độ chính xác--độ nhạy khi xem mẫu OOD là lớp dương. \\
Backbone & Mạng nền dùng để trích xuất biểu diễn đầu vào. \\
Baseline model & Mô hình đối chứng dùng để so sánh kết quả thực nghiệm. \\
Baseline input & Đầu vào tham chiếu của phép giải thích hoặc can thiệp; khác với mô hình đối chứng. \\
Bidirectional cross-attention & Chú ý chéo hai chiều: ảnh truy vấn văn bản và văn bản truy vấn ảnh. \\
Bootstrap & Lấy mẫu lại có hoàn lại để ước lượng độ bất định của một đại lượng thống kê. \\
Center crop & Cắt vùng trung tâm của ảnh sau bước thay đổi kích thước. \\
Class imbalance & Mất cân bằng lớp: số lượng quan sát chênh lệch giữa các lớp. \\
Clinical history & Bệnh sử lâm sàng; trên CTCH gồm lý do vào viện, diễn tiến bệnh và tiền sử. \\
Domain adaptation & Thích ứng miền dữ liệu; trong khóa luận là điều chỉnh biểu diễn tiền huấn luyện cho dữ liệu đích. \\
Domain shift & Dịch chuyển miền: thay đổi phân phối do cơ sở y tế, thiết bị, quần thể hoặc điều kiện thu nhận. \\
Embedding & Vector nhúng biểu diễn đầu vào trong không gian đặc trưng. \\
False-negative pair & Cặp âm giả trong học tương phản: cặp khác mẫu nhưng có quan hệ ngữ nghĩa bị xem là âm; không đồng nhất với FN phân lớp. \\
FPR@95\%TPR & Tỷ lệ chấp nhận nhầm mẫu OOD tại ngưỡng giữ lại ít nhất 95\% mẫu ID trong giao thức của khóa luận. \\
Full fine-tuning & Tinh chỉnh toàn bộ tham số mô hình. \\
Image-only & Cấu hình chỉ sử dụng ảnh, không nhận văn bản trong pha phân lớp. \\
Label smoothing & Làm trơn nhãn: phân bổ một phần xác suất đích cho các lớp khác nhằm hạn chế dự đoán quá tự tin. \\
Logit & Điểm đầu ra của mô hình trước phép Softmax. \\
Macro-F1 & Trung bình không trọng số của điểm F1 trên từng lớp. \\
Metadata & Siêu dữ liệu mô tả mẫu, như tuổi, giới tính, vị trí giải phẫu hoặc cờ nhóm bệnh. \\
Modality & Phương thức dữ liệu; khóa luận xét ảnh và văn bản. \\
Padding & Đệm chuỗi hoặc ảnh; đệm token văn bản không đồng nghĩa với đệm viền ảnh. \\
Patch token & Đơn vị biểu diễn của một mảnh ảnh trong ViT. \\
Semantic matching & Khớp ngữ nghĩa giữa biểu diễn ảnh và văn bản. \\
Shortcut learning & Học dựa vào dấu hiệu phụ tương quan với nhãn thay vì đặc trưng cần thiết cho tác vụ. \\
Token / subword & Đơn vị mã hóa văn bản / mảnh từ; token có thể là từ, mảnh từ, dấu câu hoặc ký hiệu đặc biệt. \\
Weighted cross-entropy & Entropy chéo có trọng số lớp, làm thay đổi chi phí của lỗi giữa các lớp. \\
\end{longtable}
\endgroup
